\documentclass[twocolumn]{iopjournal}

\usepackage{amsmath,amssymb,amsfonts}

\fancyhfoffset[L]{0mm}
\setlength{\marginparwidth}{0pt}
\setlength{\marginparsep}{0pt}

% ==========================================
% ZENODO PREPRINT HEADERS (ACTIVE)
% ==========================================
\renewcommand{\articletype}[1]{%
  {\vspace*{-8mm}\noindent \Large \sf Preprint}\par
  \vspace*{6mm}\noindent{\scriptsize \sf{\bfseries \MakeUppercase{#1}}}}

\fancyhead[L]{}
\fancyhead[R]{}

\begin{document}

\articletype{Letter}

\title{Gravity Field Dynamics: Field Closure of the Tetrad Frame}

\author{Stephen Nelson}

\affil{Field Dynamics Research}

\ead{stephen@fielddynamics.org}

\begin{abstract}
The companion letters derived the coupling polynomial
$f(K,s) = 1 + sK + K^2$ from the structural content of the
tetrad equation~\cite{Nelson2026b} and the constitutive law
$\mu(x) = x/(x{+}1)$ from exchange symmetry between the two
tetrahedral roots~\cite{Nelson2026a,Nelson2026b}.  This letter
proves that the torsion-free condition
$\Gamma^\lambda_{\mu\nu} = \Gamma^\lambda_{\nu\mu}$ retains
the $s = +1$ projection of the polynomial and discards the
$s = -1$ projection: the retained count per face is
$(f(K,{+}1){-}1)/2 = 10$, the discarded count is
$(f(K,{-}1){-}1)/2 = 6$, giving 40~retained and
24~discarded across four faces.  Because the companion letter
identified $s = +1$ with TetA and $s = -1$ with TetB through
the parity of the rank-1 vielbein, the 24~discarded components
are the second root of the tetrad equation.  Restoring them
yields the full $K^3 = 64$ component connection.  The
constitutive law partitions the restored components spatially:
the TetA fraction $\mu$ dominates at the galactic centre
(the bulge) and the TetB fraction $1 - \mu$ dominates at
the rim (the disk).  Integration via standard Lagrangian
mechanics yields the unique scalar-sector energy functional
$F(y) = y - 2\sqrt{y} + 2\ln(1{+}\sqrt{y})$, whose three
terms correspond to the three structural levels of the
coupling polynomial.  The constitutive law determines the
Gravity Field Dynamics covariant action; the action determines
the dynamics; no step in this chain involves a free parameter.
\end{abstract}

% ==================================================================

\section{The torsion-free projection}

\begin{figure}[!t]
\centering
\includegraphics[width=\columnwidth]{table1_matrix.png}
\caption{The $(\mu,\nu)$ matrix for one observation face
$\lambda$.  Diagonal entries (bold) are shared: the metric
cannot distinguish the two roots there.  Upper-triangle
entries belong to the $s = +1$ projection (TetA, retained).
Lower-triangle entries belong to the $s = -1$ projection
(TetB, discarded by the torsion-free condition).  The same
pattern repeats for all four values of $\lambda$.}
\label{tab:matrix}
\end{figure}

The affine connection on a $K$-dimensional frame has three
indices $\Gamma^\lambda_{\mu\nu}$, each running from 0 to
$K - 1$.  The total number of components is $K^3$.  At
$K = 4$, this gives 64~\cite{Nelson2026b,Hehl1976}.

The torsion-free condition forces symmetry of the lower pair:
\begin{equation}
  \Gamma^\lambda_{\mu\nu} = \Gamma^\lambda_{\nu\mu}.
  \label{eq:torsion_free}
\end{equation}
The companion letters established: Letter~1~\cite{Nelson2026a}
proved that the tetrad equation produces a stellated
octahedron geometry with structural constants $K = 4$,
$D = 3$, selected by the uniqueness equation
$2(D{+}1) = 2^D$.  Letter~2~\cite{Nelson2026b} showed that
the tetrad equation
$g_{\mu\nu} = \eta_{ab}\,e^a{}_\mu\,e^b{}_\nu$ contains
three structural ingredients at ranks 0, 1, and~2 on the
frame space, with graded dimensions $K^0 = 1$, $K^1 = K$,
and $K^2$.  The tetrad equation is quadratic in the frame
field and admits two roots related by spatial parity
$e^a{}_\mu \to -e^a{}_\mu$.  The metric is invariant under
this parity (both roots produce the same $g_{\mu\nu}$), so
the metric cannot distinguish the two roots.  Under parity,
a rank-$r$ tensor transforms with sign $(-1)^r$.  Ranks~0
and~2 are parity-even; rank~1 is parity-odd.  The rank-1
vielbein is therefore the only ingredient whose sign changes
when the two roots are exchanged, sending $s \to -s$ and
identifying TetA with $s = +1$ and TetB with $s = -1$.  The
coupling polynomial is
\begin{equation}
  f(K,s) = 1 + sK + K^2.
  \label{eq:poly}
\end{equation}

For each value of the upper index $\lambda$, the lower pair
$(\mu, \nu)$ forms a $K \times K$ matrix with $K^2 = 16$
entries.  Of these, $K = 4$ are diagonal ($\mu = \nu$) and
$K^2 - K = 12$ are off-diagonal.  The off-diagonal entries
form $\tfrac{1}{2}(K^2 - K) = 6$ mirror pairs, each
consisting of $(\mu, \nu)$ and $(\nu, \mu)$.  The
torsion-free condition~(\ref{eq:torsion_free}) forces the two
members of each pair to be equal, retaining the symmetric
combination and discarding the antisymmetric combination.  The
count per face is therefore:
\begin{align}
  \text{Retained:} \quad
    & K + \frac{K^2 - K}{2} = \frac{K^2 + K}{2} = 10,
    \label{eq:kept} \\[4pt]
  \text{Discarded:} \quad
    & \frac{K^2 - K}{2} = 6.
    \label{eq:disc}
\end{align}
Across $K = 4$ observation faces: $4 \times 10 = 40$ retained,
$4 \times 6 = 24$ discarded, $4 \times 16 = 64$ total.

Now observe that the coupling polynomial~(\ref{eq:poly})
evaluates to $f(K,{+}1) = 1 + K + K^2 = 21$ and
$f(K,{-}1) = 1 - K + K^2 = 13$.  Subtracting the
field-origin term ($K^0 = 1$) and halving:
\begin{align}
  \text{Retained:} \quad
    \frac{f(K,{+}1) - 1}{2}
    &= \frac{K^2 + K}{2} = 10,
    \label{eq:poly_kept} \\[6pt]
  \text{Discarded:} \quad
    \frac{f(K,{-}1) - 1}{2}
    &= \frac{K^2 - K}{2} = 6.
    \label{eq:poly_disc}
\end{align}
The combinatorial formula for the symmetric/antisymmetric
split of a $K \times K$ matrix is the coupling polynomial
evaluated at its two opposite states, with the field-origin
term removed and the mirror-pair degeneracy halved.  The
$\pm K$ that separates 10 from 6 is the $sK$ term of the
polynomial: the parity-odd rank-1 contribution.

The torsion-free condition therefore retains the $s = +1$
projection of the coupling polynomial and discards the
$s = -1$ projection.  Because the companion letter proved
that spatial parity sends $s \to -s$, identifying TetA with
$s = +1$ and TetB with $s = -1$, the retained projection is
the TetA contribution and the discarded projection is TetB.
The 24~components removed by the torsion-free condition are
the second root of the tetrad equation, carrying a distinct
field fraction $1 - \mu(x)$ at every radius.  Restoring them
returns the full $K^2 = 16$ components per face and the full
$K^3 = 64$ component connection.

The full 64-component connection is
\textit{field-closed}: both roots of the tetrad equation
contribute, the partition $\mu + (1{-}\mu) = 1$ holds at
every radius, and the field response is self-consistent
without external mass.  This is distinct from Einstein's
notion of a ``complete field theory'' (in which the equations
of motion follow from the field equations without independent
postulates) and from geodesic completeness (the extendibility
of geodesics to infinite affine parameter).  Field closure
refers specifically to the connection: whether all $K^3$
components of the affine connection are active, or whether a
subset has been set to zero by constraint.
Einstein--Cartan theory~\cite{Hehl1976} restores torsion but
couples it algebraically to spin, so that it vanishes in
vacuum.  The present framework restores torsion with an
acceleration-dependent partition: the 24~antisymmetric
components carry the fraction $1 - \mu(x)$ of the field
response, which is nonzero wherever $x$ is finite.

% ==================================================================

\section{The partition in space}

The companion letter~\cite{Nelson2026b} derived the partition
between the two projections from exchange symmetry:
\begin{align}
  \text{TetA ($s = +1$):} \quad & \mu(x) = \frac{x}{x + 1},
  \label{eq:mu} \\[4pt]
  \text{TetB ($s = -1$):} \quad & 1 - \mu(x) = \frac{1}{x + 1},
  \label{eq:tetb}
\end{align}
where $x = g/a_0$ is the dimensionless acceleration.  The
torsion-free condition forces the two projections to be equal,
but Eqs~(\ref{eq:mu})--(\ref{eq:tetb}) show they are equal
only at $x = 1$ ($\mu = 1/2$).  At every other radius, the
forced equality discards a nonzero TetB contribution: at the
galactic centre ($x \gg 1$), the discarded fraction is small;
at the rim ($x \ll 1$), it is the majority of the field
response.

What does the restored $s = -1$ projection look like
spatially?  The answer is Figure~\ref{fig1}.

\begin{figure}[t]
\centering
\includegraphics[width=\columnwidth]{fig1_tetrad_partition}
\caption{The constitutive law $\mu(x) = x/(x{+}1)$,
partitioned between the two projections identified by the
coupling polynomial.  Blue: the $s = +1$ projection $\mu$
(TetA, retained by general relativity).  Red: the $s = -1$
projection $1 - \mu$ (TetB, removed by the torsion-free
condition).  At every radius, blue~$+$~red~$= 1$.  The white
ring marks the crossover at $x = 1$, where both projections
contribute equally ($\mu = 1/2$).  The polynomial identifies
which projection is which; the constitutive law determines how
much each carries at every radius.}
\label{fig1}
\end{figure}

Figure~\ref{fig1} renders each projection as its own surface.
At each point in the radial domain, $\mu(x)$ sets the height
of the blue surface, and $1 - \mu(x)$ sets the height of the
red surface.  The two surfaces cross where their fractions are
equal, at $x = 1$.  Figure~\ref{fig1} is therefore not a
heuristic sketch but the constitutive law itself rendered over
radius: once $\mu$ and $1 - \mu$ are plotted on the same
domain, the bulge-to-disk ordering follows directly.

Every gravitationally bound system has an acceleration
profile that is strongest at the centre and weakens outward.
This is not a model; it is Gauss's law applied to any finite
mass distribution.  The constitutive law
$\mu(x) = x/(x{+}1)$ is a monotonic function: higher
acceleration always maps to a larger TetA fraction, and lower
acceleration always maps to a larger TetB fraction.  The
composition of a monotonic partition with a monotonically
decreasing acceleration profile produces the same
bulge-to-disk radial ordering.  Where $x$ is large (the
centre), TetA dominates and the blue surface peaks.  Where
$x$ is small (the rim), TetB dominates and the red surface
rises.  The result is a peaked central concentration inside
an extended disk --- a bulge inside a disk --- for any
gravitationally bound system, not only for a particular
galaxy model.  The resemblance between Figure~\ref{fig1} and
a spiral galaxy is therefore not a visual coincidence but the
direct rendering of the constitutive law over radius.

The vertical axis in Figure~\ref{fig1} is the field fraction,
not a physical height.  The blue spike at centre is TetA
carrying nearly the entire field response.  The red plateau at
large radius is TetB carrying nearly all of it.  Between them,
the white ring marks the handover.  The textbook decomposition
of the full connection into Christoffel plus contorsion
corresponds, at the level of field magnitudes, to the
partition into $\mu$ and
$1 - \mu$~\cite{Nelson2026b}.  Figure~\ref{fig1} is that
decomposition rendered spatially.

The white ring sits at $x = 1$, where $\mu = 1/2$.  This is
the locus where both projections contribute equally to the
field response.  Inside the ring, the Christoffel sector
dominates: general relativity is accurate and Keplerian
decline holds.  Outside the ring, the torsion sector
dominates: the field response that general relativity
discards becomes the majority.  This crossover is the MOND
transition scale $a_0$~\cite{Milgrom1983}.  It is not a
fitted parameter.  It is the geometric midpoint of the
exchange-symmetric partition, the point where the odds ratio
$\Omega = \mu/(1-\mu)$ equals unity.

The rate at which the partition shifts between the two
projections is characterised by a logistic identity that
$\mu$ satisfies~\cite{Nelson2026b}:
\begin{equation}
  x\,\frac{d\mu}{dx} = \mu(1 - \mu).
  \label{eq:logistic}
\end{equation}
Under the substitution $t = \ln x$, this becomes
$d\mu/dt = \mu(1-\mu)$: the canonical two-sector transition
law.  The product $\mu(1-\mu)$ on the right-hand side measures
the simultaneous presence of both tetrahedral chambers at a
given radius.  It is the continuous analogue of the octahedral
interaction core, the $s = 0$ state of the coupling
polynomial~\cite{Nelson2026b}.  The transition peaks at the
MOND ring because that is where the core is fully realised:
$\mu(1 - \mu) = 1/4$.

The kinematic consequence follows directly from the partition.
At dimensionless acceleration $x = g_{\mathrm{obs}}/a_0$, the
baryonic acceleration, carried by the Christoffel sector, is
$g_{\mathrm{bar}} = \mu\,g_{\mathrm{obs}}$, giving
\begin{equation}
  g_{\mathrm{bar}} =
  \frac{g_{\mathrm{obs}}^2}{g_{\mathrm{obs}} + a_0}.
  \label{eq:rar}
\end{equation}
At high acceleration ($x \gg 1$):
$g_{\mathrm{bar}} \approx g_{\mathrm{obs}}$, Keplerian
decline, general relativity is accurate.  At low acceleration
($x \ll 1$):
$g_{\mathrm{bar}} \approx g_{\mathrm{obs}}^2/a_0$, flat
rotation curves.  The radial acceleration relation established
empirically by McGaugh, Lelli \&
Schombert~\cite{McGaugh2016}, with 0.057~dex scatter across
175~individual galaxies~\cite{Li2018} and no credible
variation in $a_0$, is reproduced by Eq.~(\ref{eq:rar}) with
no free parameters.  The full 179-galaxy rotation curve
analysis is presented in a companion
paper~\cite{Nelson2026c}.

The apparent separation of baryonic and dark components in
Figure~\ref{fig1} is a diagnostic decomposition of one field,
not a claim of two independent physical substances.  It is the
projection of a single 64-component field onto a 40-component
subspace.  The observed gravitational discrepancy at galactic
scales is not in dispute: both the cold dark matter framework
and the present partition reproduce the radial acceleration
relation and the function $\mu(x) = x/(x{+}1)$.  They differ
on what the discrepancy represents.  The cold dark matter
framework interprets it as an invisible mass component,
modelled with approximately spherical
halos~\cite{NFW1996}.  The closed tetrad frame identifies it
as the field response of the second tetrahedral chamber,
carried by the 24~torsion components.  This response is not
mass.  It is topology.

% ==================================================================

\section{From partition to action}

The partition developed in Sections~1--2 produces three
observational consequences that distinguish it from the cold
dark matter framework.

\textit{Disk-and-bulge morphology.}  Because the constitutive
law modulates the torsion fraction radially, the partition
produces a Christoffel-dominated bulge at the centre
($\mu \approx 1$) transitioning into a torsion-dominated disk
at the rim ($\mu \approx 0$), with the morphological
transition at the MOND ring $x = 1$.

\textit{Bounded central density.}  The TetB field fraction
$1/(x{+}1)$ is bounded above by unity for all $x$.  The
torsion-sector field acceleration never exceeds $a_0$.  The NFW
profile diverges as $r \to 0$~\cite{NFW1996}.  The partition's
boundedness removes the central divergence responsible for the
cusp-core tension, without baryonic
feedback or self-interacting dark matter.  The removal is
algebraic: $1/(x{+}1) < 1$ is a property of the constitutive
law, not of any galaxy model.

\textit{No new particles.}  The gravitational discrepancy is
not mass.  It is the field response of the second tetrahedral
chamber, carried by the 24~torsion components that the
zero-torsion condition removes from the connection.  The
missing component is not in the matter sector.  It is in the
connection.

General relativity succeeds in the solar system because the
solar system sits deep inside a galaxy where the gravitational
acceleration exceeds the transition scale $a_0$ by a factor of
order $10^8$.  At those accelerations the partition fraction
$\mu$ is so close to unity that the second root contributes
less than one part in a million.  The torsion-free
approximation loses nothing measurable, and general relativity
was confirmed on every local test it was given.  But a galaxy
spans the full range of the partition: from the core where one
root carries nearly everything, through the crossover ring
where both roots contribute equally, to the rim where the
second root carries the majority.  The function
$\mu(x) = x/(x{+}1)$ connects these regimes smoothly.
Milgrom~\cite{Milgrom1983} identified it empirically.  Famaey
\& McGaugh~\cite{FamaeyMcGaugh2012} selected it by fitting
galactic data.  The constitutive law that appears
phenomenologically in the AQUAL formulation of modified
gravity is here recovered as the geometric partition of the
two-chamber tetrad frame.  The present series derives it: the
tetrad equation is degree two, exchange symmetry between its
roots admits exactly one rational partition, and that
partition is $\mu$.

The constitutive law also determines the scalar-sector
energy functional.  The variational relationship between a
constitutive law and its energy functional is standard
Lagrangian mechanics: the response function is the derivative
of the potential.  Applied to gravity, this is the identity
$\mu(\sqrt{y}) = dF/dy$ used by Bekenstein \&
Milgrom~\cite{BekensteinMilgrom1984} in the AQUAL
formulation, where
$y = |\nabla\Phi|^2/a_0^2 = x^2$ is the squared
dimensionless acceleration.  Integration of
$\mu(x) = x/(x{+}1)$ yields the unique kinetic functional
\begin{equation}
  F(y) = y - 2\sqrt{y} \;+\; 2\ln\!\bigl(1 + \sqrt{y}\bigr),
  \label{eq:F}
\end{equation}
In the standard AQUAL literature, $F$ is a free function
selected to fit observations.  Here it is fixed by the
constitutive law, which is itself fixed by the frame geometry.
The three terms of Eq.~(\ref{eq:F}) correspond to the three
structural levels of
the coupling polynomial~\cite{Nelson2026b}: the quadratic $y$
to the $K^2$ face-coupling level, the $-2\sqrt{y}$ to the
state-dependent $sK$ propagation level, and the logarithm to
the $K^0$ field origin.  At high acceleration ($y \gg 1$),
$F \to y$ and the action reduces to general relativity.  At
low acceleration ($y \ll 1$), $F \to \tfrac{2}{3}\,y^{3/2}$,
recovering the deep-field regime and the baryonic Tully-Fisher
relation.  The constitutive law determines the action.  The
action determines the dynamics.  No step in this chain
involves a free parameter.

The functional $F(y)$ determines the scalar-sector kinetic
term of the Gravity Field Dynamics (GFD) covariant action:
\begin{align}
  S_{\text{ST}} &= \int\!d^4x\,\sqrt{-g}\biggl[
    \frac{R}{16\pi G}
    - \frac{a_0^2}{8\pi G}\,F\!\bigl(|\nabla\Phi|^2/a_0^2\bigr)
    \notag\\
    &\qquad\qquad
    -\; \tfrac{1}{4}\,e^{-2\Phi/c^2}
    F_{\mu\nu}F^{\mu\nu}
    + \mathcal{L}_{\text{matter}}
  \biggr].
  \label{eq:sst}
\end{align}
The first term is the Einstein--Hilbert action of general
relativity.  The second is the scalar sector, whose kinetic
function $F$ is derived above from the constitutive law.  The
third is a disformal electromagnetic coupling.  The fourth is
standard matter.  The derivation of the full action from the
constitutive law, including the electromagnetic coupling and
its consequences for gravitational lensing, is provided in
the supplementary covariant completion note.  No term
contains a free function: the constitutive law fixes $F$, and
the frame geometry determines the remaining couplings.

To summarise the downstream structure: the partition
$\mu = x/(x{+}1)$ is selected uniquely by exchange symmetry
between the two tetrad roots~\cite{Nelson2026b}.  The same
function partitions the 64-component connection into the
40~Christoffel and 24~torsion sectors.  Its complement
$1 - \mu$ is exactly the field fraction that, in the
torsion-free projection, must be attributed to invisible mass.
It satisfies the logistic transition equation
$d\mu/d\!\ln x = \mu(1{-}\mu)$.  Its integral yields the
unique AQUAL energy functional $F(y)$, whose three terms map
to the three structural levels of the coupling polynomial.
That functional recovers general relativity at high
acceleration and the baryonic Tully-Fisher relation at low
acceleration.  These are not separate results.  They are the
same algebraic object encountered at successive levels of the
theory: partition, connection, transition, action, kinematics.
These closures are not independent.  Rejection of any one
therefore requires an alternative account of the whole chain.

An alternative relativistic approach to the MOND phenomenology
has been developed by Skordis \&
Z\l{}o\'{s}nik~\cite{SkordisZlosnik2021}, who introduce
additional vector and scalar fields to obtain a theory
compatible with the cosmic microwave background power spectrum.
The present framework adds no fields.  It restores the
24~connection components that general relativity's zero-torsion
condition removes, recovering the complete geometry of the
tetrad frame.  The result is field closure: the gravitational
connection carries all $K^3 = 64$ components, the two
tetrahedral chambers contribute their full partition at every
radius, and the field response balances without invisible mass.

The torsion-free specialisation adopted in 1915 discards the
24~antisymmetric components of the affine connection.
Restoring them closes the field.  The field response they
carry is precisely what, in their absence, must be attributed
to invisible mass.

% ==================================================================

\section{Data and reproducibility}

An interactive rendering of Figure~\ref{fig1} is provided as
a supplementary HTML file permitting rotation, zoom, and
background toggling.  A self-contained verification script
(\texttt{validate\_letter3.py}) confirms the polynomial
identity~(\ref{eq:poly_kept})--(\ref{eq:poly_disc}), the partition
properties~(\ref{eq:mu})--(\ref{eq:tetb}), the logistic
identity~(\ref{eq:logistic}), the radial acceleration
relation~(\ref{eq:rar}), and both asymptotic limits of
$F(y)$.  To complete the field-closure claim, a supplementary
note (\texttt{supplementary\_covariant\_completion.pdf})
derives the full four-term covariant action from the
constitutive law, confirming that the closed connection yields
a complete relativistic field theory with no free functions.

% ==================================================================

\begin{thebibliography}{99}

\bibitem{Nelson2026a}
S.~Nelson,
\textit{The Coordinate-Space Geometry of the Tetrad Frame},
Class.\ Quantum Grav.\ (2026), submitted (CQG-115428).

\bibitem{Nelson2026b}
S.~Nelson,
\textit{The Other Half of Gravity: Torsion and the
Two-Chamber Structure of the Tetrad Frame},
Class.\ Quantum Grav.\ (2026), submitted (CQG-115447).

\bibitem{Nelson2026c}
S.~Nelson,
\textit{Gravity Field Dynamics: Decoding mass from observed
rotation curves},
Mon.\ Not.\ R.\ Astron.\ Soc.\ (2026), submitted.

\bibitem{McGaugh2016}
S.~S.~McGaugh, F.~Lelli and J.~M.~Schombert,
\textit{Radial Acceleration Relation in Rotationally
Supported Galaxies},
Phys.\ Rev.\ Lett.\ \textbf{117} (2016) 201101.

\bibitem{Li2018}
P.~Li, F.~Lelli, S.~McGaugh and J.~Schombert,
\textit{Fitting the radial acceleration relation to
individual SPARC galaxies},
Astron.\ Astrophys.\ \textbf{615} (2018) A3.

\bibitem{Milgrom1983}
M.~Milgrom,
\textit{A modification of the Newtonian dynamics as a
possible alternative to the hidden mass hypothesis},
Astrophys.\ J.\ \textbf{270} (1983) 365.

\bibitem{FamaeyMcGaugh2012}
B.~Famaey and S.~S.~McGaugh,
\textit{Modified Newtonian Dynamics (MOND): Observational
Phenomenology and Relativistic Extensions},
Living Rev.\ Relativity \textbf{15} (2012) 10.

\bibitem{BekensteinMilgrom1984}
J.~Bekenstein and M.~Milgrom,
\textit{Does the missing mass problem signal the breakdown
of Newtonian gravity?},
Astrophys.\ J.\ \textbf{286} (1984) 7.

\bibitem{Hehl1976}
F.~W.~Hehl, P.~von~der~Heyde, G.~D.~Kerlick and
J.~M.~Nester,
\textit{General relativity with spin and torsion:
Foundations and prospects},
Rev.\ Mod.\ Phys.\ \textbf{48} (1976) 393.

\bibitem{NFW1996}
J.~F.~Navarro, C.~S.~Frenk and S.~D.~M.~White,
\textit{The structure of cold dark matter halos},
Astrophys.\ J.\ \textbf{462} (1996) 563.

\bibitem{SkordisZlosnik2021}
C.~Skordis and T.~Z\l{}o\'{s}nik,
\textit{New Relativistic Theory for Modified Newtonian
Dynamics},
Phys.\ Rev.\ Lett.\ \textbf{127} (2021) 161302.

\end{thebibliography}

\end{document}
